\section{Methods and Experimental Setup}
\label{sec:experimental_setup}

\subsection{Architectures and Benchmark}
\label{sec:setup_overview}

Seven VLA architectures spanning three vision connector families provide the testbed (Table~\ref{tab:taxonomy}).

\textbf{Linear projection.} $\pi_0$~\cite{black2024pi0} (3.3B, 256 tokens) maps SigLIP~\cite{zhai2023siglip} patches through a single linear layer; $\pi_0.5$~\cite{black2025pi05} (4B, 768 tokens from three cameras) adds a learned $\times\!\sqrt{d}$ scaling; X-VLA~\cite{zheng2025xvla} (880M, 50 tokens) augments its linear projection with a 12-layer BART encoder~\cite{lewis2020bart} that introduces cross-token mixing. \textbf{MLP projector.} OpenVLA~\cite{kim2024openvla} (7B, 256 tokens) fuses DINOv2~\cite{oquab2023dinov2} and SigLIP through a three-layer MLP; OFT~\cite{kim2025openoft} shares this connector (entropy comparison only); CogACT~\cite{li2024cogact} retains 74.2\% SR under 78\% pruning per EfficientVLA~\cite{yang2025efficientvla} (literature corroboration, not our evaluation). \textbf{Spatial folding.} SmolVLA~\cite{shukor2025smolvla} (450M, 128 tokens) applies pixel-shuffle ($4{\times}4$ folding) so each token encodes 16 contiguous patches.

We evaluate on LIBERO~\cite{liu2023libero} (Object: 10 tasks $\times$ 50 episodes = 500 total; Spatial and Goal for cross-suite validation). Closed-loop success rate (SR) with Wilson 95\% CIs is the primary metric; multi-seed results report mean $\pm$ std. Baselines: SmolVLA 44.8\%, OpenVLA 74.0\%, X-VLA 99.2\%, $\pi_0.5$ 98.2\%.

\subsection{Pruning Protocol}
\label{sec:pruning_protocol}

Two methods bracket robustness. \textbf{L2 norm-based removal} ranks tokens by L2 norm and physically removes the lowest-norm fraction, providing a favorable test. \textbf{Random removal} selects tokens uniformly at random, providing a conservative lower bound. Both physically remove tokens (not zero-masking); remaining tokens retain spatial ordering with re-indexed position IDs. Per-model retention ratios and cross-architecture comparison details appear in Appendix~\ref{app:pruning_details}.

\subsection{Vision Token Re-encoder (VTR)}
\label{sec:smart_compress_setup}

A cross-attention module (22.2M parameters, $\sim$5\% of SmolVLA) is inserted between the pixel-shuffle connector and the LLM backbone. $M$ learned queries attend over 128 source tokens, producing $M$ compressed representations ($M \in \{16, 32, 64, 96, 128\}$; $M{=}128$ = passthrough control). A parameter-matched MLP adapter (22.14M) serves as architecture control. Both train from SmolVLA-Base for 100K steps with identical hyperparameters; an extended training control without the compressor achieves 45.0\% SR ($p > 0.9$ vs.\ 44.8\% baseline). Training details in Appendix~\ref{app:smart_compress_details}.

\subsection{Diagnostic Metrics}
\label{sec:diagnostic_metrics}

\textbf{Token importance CV} (coefficient of variation of L2 norms) tracks pruning tolerance continuously: uniformly critical at CV $< 0.06$, threshold-graded at $\sim$0.15, smooth graded at $\sim$0.22, catastrophic cliff at $> 0.3$. \textbf{Attention entropy ratio} (normalized Shannon entropy of attention weights over vision tokens) provides a complementary view: OFT and OpenVLA share nearly identical entropy (0.667 vs.\ 0.661), confirming the connector is associated with token utilization patterns independent of action head.
